%%LTPROPv0.7     DO NOT REMOVE/EDIT THIS LINE!                
%%
%%      This is v0.7 from 1st June 2017 (RJS)  
%%      This for is an adaptation of previous LT PATT forms, which ultimately 
%%      trace back to a JCMT version by gms@jach.hawaii.edu 
%%
%%      The current style file (ltpattv0p7) must be installed for 
%%      this template to work.   
%%
%%      Proposals submitted based on older versions of this template or
%%      the older style file will be rejected by the proposal submission
%%      system.  To check you have the latest version, see
%%
%%      http://telescope.livjm.ac.uk/
%%
%%      The above web site also has useful details to help you prepare
%%      your proposal
%%
%%       For help and support please email ltphase1@ljmu.ac.uk
%%
%% To use this file simply fill in the gaps between the curly brackets
%% with the appropriate info.  DO NOT fiddle with anything else as this
%% will lead to the automated submission system rejecting your proposal.
%% Similarly DO NOT attempt to edit the style file.
%%
%% If you prefer to use psfig rather than epsf to embed your figures 
%% comment out the first \usepackage and uncomment the second one.
%% 

\documentclass[11pt]{article}
\usepackage{ltpattv0p7}
%\usepackage{psfig}
%\usepackage{epsf}
%\usepackage{amssymb}

\begin{document}

\tag{} % PATT or JMU

\semester{} % semester of proposal (e.g. 01A, 01B, 02A etc.)
\duration{}   % how many semesters does this proposal cover?

\progid{} % If the proposal is part of an ongoing LT programme, enter the
               % PROGID here.  

\PIname  {}        % PI's name
\PIinstitution {}  % institution, and country
\PIEMAIL {}        % IF YOUR EMAIL ADDRESS INCLUDES
                   % UNDERSCORE CHARACTERS, YOU MUST
                   % PRECEDE THEM BY A BACKSLASH
                   % E.G. fred\_bloggs@univ.ac.uk
                                  
 
% you can have up to 5 Co-I's - they don't need EMAIL's, but if they
% don't have them they will not be able to participate in phase 2 data
% entry or download data from JMU.

\CIAname  {}  
\CIAinstitution {}
\CIAEMAIL {}

\CIBname  {}  
\CIBinstitution {}
\CIBEMAIL {}

\CICname  {}  
\CICinstitution {}
\CICEMAIL {}

\CIDname  {}  
\CIDinstitution {}
\CIDEMAIL {}

\CIEname  {}  
\CIEinstitution {}
\CIEEMAIL {}

% Please type Y in the element below if this application is for data to be used in a PhD thesis.
% Also enter the name and institute for the PhD candidate

\phdthesis{}
\phdname{}

% The proposal title should fit on one line of the final form.

\title{}

% Enter the abstract below.  Do not change the font size.

\abstract{}

%Enter the amount of time you are requesting THIS SEMESTER.
%The unit of time is HOURS.

\thissemestertime{}

% If you are applying for a programme length of > 1 semester then
% enter the amount of time you are requesting IN TOTAL.
% The unit of time is HOURS.

\totaltime{}

% Enter the required seeing value in arcseconds.

\seeing{}

% Enter the required sky brightness in units of magnitudes per
% square arcsecond above DARK.  See the telescope website
% for details.

\sky{}

% Enter Y below if photometric conditions required, N otherwise

\photometric{}

% Enter Y in the months in which you would ideally like your observations
% to take place.  you will get an opportunity to put more detailed
% constraints on your actual observations at phase 2 - this is just to 
% allow a balanced (in terms of season) set of allocations to be made 
% by the TAG.

\jan{}
\feb{}
\mar{}
\apr{}
\may{}
\jun{}
\jul{}
\aug{}
\sep{}
\oct{}
\nov{}
\dec{}


%If observation is a Target of Opportunity (i.e. you will be responsible for
%entering the phase 2 information when an event occurs, rather than at the
%start of the semester) please enter Y here.  Otherwise enter N

\too{}

%If observation is a Target of Opportunity enter an per-cent estimate of the
%likelihood of the trigger occurring during a single semester.

\likelihood{}  %DO NOT USE A % SIGN IN FILLING THIS IN!  JUST GIVE THE FIGURE.

%Enter details of timing constraints for groups. Options are any
%combination of FLEXIBLE, FIXED, MONITOR, INTERVAL, PHASED

\constraints{}

%If you have a timing constraint that implies repeated observations (not simply
%multiple exposures in a single observation), please indicate the mean time
%between observations below in DAYS

\cadence{}

%Indicate the minimum useable fraction (as a per-centage) of observations 
%necessary to obtain your science objectives (remember that for
%example a monitoring campaign is likely to be affected
%by bad weather, so it is very unlikely to achieve 100%).

\muf{} %DO NOT USE A % SIGN IN FILLING THIS IN!  JUST GIVE THE FIGURE.

% To use the IO:O CCD Camera, enter Y in the box below.
\instioo{}
% If using IO:O, please enter details of your filters below.
\instioodetails{}

%to use the the IO:I near-IR camera, enter Y in the \instioi{} box
%below. 
\instioi{}
% If using IO:I, please enter details of your filters below.
\instioidetails{}

%to use the RINGO3 polarimeter, enter Y in the \instringo box below.
\instringo{}

%to use the RISE fast-photometer, enter Y in the \instrise box below.
\instrise{}

%to use the SPRAT long-slit spectrometer, enter Y in the \instsprat{} box
%below.  
\instsprat{}

%to use the FRODOSPEC spectrographs (depending on
%enter Y in the relevant box below.
\instfrodo{}
%Enter details of FRODOSPEC gratings below
\instfrododetails{}

%to use another instrument, enter its name in the \instother box below.
%you MUST consult with JMU at least two weeks in advance of submission
%to do this.  JMU reserve the right not to host any particular
%instrument.
\instother{}

% Enter a summary of progress on previous LT proposals that are
% in some way related to this proposal. 
%
% Use the table format provided to report past allocations in hours, 
% the A,B,C grade allocated by the TAG and how much of the allocation was
% used. Use the ?Comments? column for any extra information that may be 
% pertinent.
%
% There is also space after the table for any text commentary you
% wish to include.
\previouswork{

\hfil\break
\begin{tabular*}{0.9\textwidth}{lrrrp{0.53\textwidth}}
\hline
{\bf PROPID} & {\bf Alloc.} & {\bf Rank} & {\bf Used} & {\bf Comments}\\
\hline
%Five columns: PROPID, hours allocated, rank, hours used, free-form comments
PL00A00 & 1 & A & 1 & {\em Delete this table line and add your own data}\\
PL00B00 & 1 & B & 1 & {\em Delete this table line and add your own data}\\
%Some examples for guidance
%PL10A35 & 5 & B & 1.5 & Low usage resulted from lack of suitable triggers and then
%an extended period of bad weather that clashed with the one good ToO trigger we got.\\
%PL10B34 & 4 & B & 4 & Successful. See publication list\\
\hline
\end{tabular*}

{\em Delete this sentence and enter any further information you wish, beyond the table.}

}

%Enter a list of all publications either published or in press
%(but not "in preparation") which are wholly or partly based on
%LT data for the past 2 years.  This should include publications
%both related and unrelated to this proposal.

\publications{

}

% Enter details of related/complimentary proposals below on
% both PATT and NON-PATT facilities, mentioning
% any resulting scheduling constraints that we should be aware of.

\otherfacilities{}

% Enter details of your targets below, or for large samples, details of the
% selection and total number.  Coordinates etc. need not be exact
% at this stage as you will do all this in more detail at phase 2 -
% this is just to allow the TAG to judge your science case.

\targetlist{

{\em Delete these italicised instructions and enter your targets below.
For large samples, give details of the target selection and total number.
Coordinates etc. need not be exact at this stage as you will do all this 
in more detail at Phase 2. This is just to allow the TAG to judge your science case.}
}

% Enter your science case between the curly brackets below.  Do NOT
% fiddle with the font size - your application will be automatically
% rejected if you do so.

\sciencecase{

{\em Delete all these italicised instructions and enter your scientific case here.

Do not exceed one page.
 
Requests for a time allocation over multiple semesters should be specifically justified here. 

If the time request is for a target-of-opportunity, the likelihood or rate of the event(s) should also be discussed. }

}

%enter the technical case below, including details of 
%signal to noise requiered and expected and why you have
%applied for the lunar and atmospheric condtions chosen.  Note that SNR
%calculators are available on the web page (http://telescope.livjm.ac.uk/)

\technicalcase{

{\em Delete all these italicised instructions and enter your technical case here.

Do not exceed one page.
 
Those requiring FIXED timing constraints for their
observations must justify this here. See the LT Phase 1 webpage
for further details.

Users who require more than a one-year post
end-of-semester proprietary period on their data must likewise
justify this here. The LT's policy on data access and proprietary
periods is posted on the LT Data Products webpage.  }

}

%enter any tabular data, references, and include any postscript
%figures using psfig or epsf on the page below.

\figsandtabs{

{\em Delete this sentence and enter figures and tables here: 
do not exceed one page.}

% to include an EPS figure use something like the line below...
% (make sure the bounding box is sensible!)
%   \psfig{figure=figname.ps,height=8cm} 

}


%	END OF INPUT  --  DON'T FIDDLE WITH ANYTHING BELOW THIS POINT
%_____________________________________________________________________________
%
%	
%	The following actually produces the pages.
%
 
 \frontpage  
 \pagebreak
 \frontpagepart2  
 \pagebreak
 \secondpage 
 \pagebreak
 \thirdpage  
 \pagebreak
 \fourthpage 
 \pagebreak
 \fifthpage  
 \pagebreak

\end{document}

